\documentclass[../../../../master.tex]{subfiles}
\begin{document}

\monographleaf{Domain 07: Agentic AI \& Production ML Systems}{Agentic Reasoning, Memory, RAG, MCP/A2A \& Workflows}{RAG, Context Engineering, Memory \& Tool/Agent Protocols}{Function Calling, Model Context Protocol (MCP), A2A Interoperability \& Workflow Automation}{Agent Tools \& MCP Whitepaper, MCP Illustrated Guidebook, Prototype to Production, RLHF Ch.~12, n8n Guide}

% =============================================================================
\section{Problem Formulation \& Unified Notation}
\label{sec:mcp:formulation}
% =============================================================================

Modern agentic systems transcend the static knowledge cutoff and probabilistic arithmetic of pretrained weights by coupling autoregressive token generation with external execution environments. Given a user query $q$ and a set of $N$ external capabilities $\mathcal{T} = \{t_1, \dots, t_N\}$ exposed across $M$ AI host applications $\mathcal{H} = \{h_1, \dots, h_M\}$, an agentic runtime must solve two coupled problems: (1) \emph{syntactic and semantic grounding}, where a language model emits schema-conformant invocations interleaved with reasoning tokens \bookcite{reinforcement_learning_from_human_feedback}{Ch.~13, PDF pp.~167--175}, and (2) \emph{protocol-level interoperability}, where $M$ heterogeneous hosts and $N$ heterogeneous tools, data sources, and peer agents communicate without incurring an $\mathcal{O}(M \cdot N)$ bespoke integration explosion \bookcite{agent_tools_interoperability_mcp}{PDF p.~20}.

To unify training-time function calling, the Model Context Protocol (MCP), the Agent2Agent (A2A) protocol, and deterministic workflow DAGs (such as n8n), we model an agentic execution trajectory $\tau$ as an alternating sequence of model-emitted actions and environment observations:
\begin{equation}
  \tau = \bigl(x_0, a_1, o_1, a_2, o_2, \dots, a_K, o_K, y_{\text{final}}\bigr),
  \label{eq:mcp:trajectory}
\end{equation}
where $x_0$ encodes the system instructions, tool/agent manifests, and user prompt; $a_k \in \mathcal{A}_{\text{tool}} \cup \mathcal{A}_{\text{a2a}}$ is a structured invocation (either an MCP JSON-RPC tool call or an A2A task delegation); $o_k$ is the returned structured/unstructured observation (or error payload); and $y_{\text{final}}$ is the terminal synthesis. When embedded inside a deterministic workflow graph $G = (V, E)$, each node $v \in V$ is either a deterministic operator $f_v$ (e.g., webhook trigger, HTTP request, filter, or schema validator) or a bounded stochastic agentic sub-trajectory $\tau_v$ \bookcite{ultimate_n8n_starter_kit_2025}{Modules~2--4, PDF pp.~8--31}.

\begin{crossbook}[Cross-Book Notation \& Pedagogical Reconciliation]
\label{tab:mcp:notation}
Across the five mapped sources, overlapping terminology describes distinct layers of the tool-calling and agent-interoperability stack. Table~\ref{tab:mcp:reconcile} reconciles these definitions into a unified architectural vocabulary.

\vspace{0.4em}
\noindent\begin{tabularx}{\linewidth}{@{} >{\raggedright\arraybackslash}p{2.3cm} >{\raggedright\arraybackslash}X >{\raggedright\arraybackslash}p{3.5cm} @{}}
\toprule
\textbf{Concept} & \textbf{Cross-Book Usage \& Scope} & \textbf{Unified Monograph Term} \\
\midrule
Tool Use vs.\ Function Calling &
  \bookcite{reinforcement_learning_from_human_feedback}{PDF p.~167} distinguishes general \emph{tool use} (orchestrator executes model request) from \emph{function calling} (arguments strictly conform to JSON Schema). &
  Schema-constrained function calling ($\mathcal{A}_{\text{tool}}$) \\
\addlinespace
MCP Host / Client / Server &
  \bookcite{agent_tools_interoperability_mcp}{PDF pp.~21--22} and \bookcite{mcp_2025_editionthe_illustrated_guidebook}{PDF pp.~10--12} define the Host (user-facing app), Client (1:1 session adapter inside Host), and Server (capability wrapper). &
  Three-tier MCP vertical stack ($\mathcal{H} \to \mathcal{C} \to \mathcal{S}$) \\
\addlinespace
MCP Primitives &
  \bookcite{mcp_2025_editionthe_illustrated_guidebook}{PDF pp.~13--18} categorizes control: Tools (model-controlled), Resources (app-controlled, read-only URIs), Prompts (user-controlled templates). \bookcite{agent_tools_interoperability_mcp}{PDF pp.~24--25} adds client-side Sampling, Elicitation, and Roots. &
  Server primitives ($\mathcal{P}_{\text{srv}}$) vs.\ Client primitives ($\mathcal{P}_{\text{cli}}$) \\
\addlinespace
MCP vs.\ A2A &
  \bookcite{prototype_production}{PDF pp.~27, 34} contrasts MCP (``Do this specific thing''---stateless/structured tool access) with A2A (``Achieve this complex goal''---stateful, autonomous agent collaboration via Agent Cards). &
  Vertical tool plumbing (MCP) vs.\ Horizontal agent delegation (A2A) \\
\addlinespace
Agentic vs.\ Deterministic Nodes &
  \bookcite{ultimate_n8n_starter_kit_2025}{PDF pp.~8--11, 26--31} separates deterministic Trigger/Core nodes (IF, Merge, Code, Set) from stochastic AI Agent nodes requiring post-processing validation. &
  Hybrid DAG $G=(V_{\text{det}} \cup V_{\text{agent}}, E)$ \\
\bottomrule
\end{tabularx}
\label{tab:mcp:reconcile}
\end{crossbook}

% =============================================================================
\section{Core Mathematical Derivations \& Theorems}
\label{sec:mcp:theory}
% =============================================================================

\subsection{The \texorpdfstring{$M \times N$}{M x N} Integration Bottleneck \& Protocol Hub Reduction}

Prior to standardized agent protocols, \bookcite{agent_tools_interoperability_mcp}{PDF p.~20} and \bookcite{mcp_2025_editionthe_illustrated_guidebook}{PDF pp.~8--9} note that connecting $M$ AI host applications to $N$ external tools required custom point-to-point adapters for every $(h_i, t_j)$ pair. Introducing an intermediate protocol hub (MCP for tools/resources, A2A for peer agents) decouples host implementations from server implementations.

\begin{coremath}[Connector Complexity Theorem: Bilateral Glue vs.\ Standardized Hub]
Let $\mathcal{H} = \{h_1, \dots, h_M\}$ be the set of AI host applications and $\mathcal{T} = \{t_1, \dots, t_N\}$ be the set of external tools or services.
\begin{enumerate}
  \item \textbf{Point-to-Point Regime:} Without a shared wire protocol, each pair $(h_i, t_j) \in \mathcal{H} \times \mathcal{T}$ requires a bespoke connector $c_{ij}$ mapping $h_i$'s internal function-calling format to $t_j$'s API schema, yielding total integration maintenance cost:
  \begin{equation}
    C_{\text{bespoke}}(M, N) = \kappa_{\text{glue}} \cdot M \cdot N = \mathcal{O}(M N).
    \label{eq:mcp:mn_cost}
  \end{equation}
  \item \textbf{Standardized Hub Regime (MCP / A2A):} Each host $h_i$ implements a single protocol client adapter $c^{\text{cli}}_i$ (speaking JSON-RPC~2.0), and each tool provider $t_j$ implements a single protocol server $s^{\text{srv}}_j$ exposing JSON Schema manifests. The total connector count reduces to:
  \begin{equation}
    C_{\text{proto}}(M, N) = \kappa_{\text{cli}} M + \kappa_{\text{srv}} N = \mathcal{O}(M + N),
    \label{eq:mcp:m_plus_n_cost}
  \end{equation}
  achieving an asymptotic reduction factor $\rho(M, N) = \frac{MN}{M + N}$, which grows as $\frac{N}{2}$ when $M = N$.
\end{enumerate}
\end{coremath}

\begin{intuition}[Why LSP Inspired MCP and Why A2A Completes the Stack]
Just as the Language Server Protocol (LSP) allowed $M$ code editors to support $N$ programming languages via $M + N$ implementations, MCP standardizes the \emph{vertical} boundary between an AI Host and its deterministic Tools/Resources \bookcite{agent_tools_interoperability_mcp}{PDF p.~21}. However, when an enterprise deploys multiple specialized agents built across heterogeneous frameworks (e.g., CrewAI, LangGraph, ADK), treating an autonomous peer agent as a stateless MCP tool breaks down: peer agents maintain multi-turn state, negotiate plans, and execute long-running tasks. A2A provides the complementary \emph{horizontal} protocol for agent-to-agent delegation \bookcite{prototype_production}{PDF pp.~27--34}.
\end{intuition}

\subsection{Constrained Decoding, Interleaved Tool Trajectories \& Post-Training}

At inference time, a language model $\pi_\theta$ generates tokens autoregressively. During a tool-augmented trajectory (Eq.~\ref{eq:mcp:trajectory}), the orchestrator pauses decoding when $\pi_\theta$ emits a tool-call delimiter or structured JSON block, executes the call, appends $o_k$ to the context, and resumes generation \bookcite{reinforcement_learning_from_human_feedback}{Ch.~13, PDF pp.~167--174}.

\begin{coremath}[Schema-Constrained Decoding \& Policy Optimization over Tool Trajectories]
Let $\mathcal{G}_{\text{JSON}}(\sigma_t)$ denote the context-free grammar induced by a tool's JSON Schema $\sigma_t = (\text{inputSchema}, \text{required})$. Under constrained decoding (e.g., finite-state masking in vLLM or provider ``strict mode''), at decoding step $s$ with prefix $w_{<s}$, the admissible next-token set is restricted to $\mathcal{V}(w_{<s}) = \{v \in \mathcal{V} : w_{<s} \cdot v \text{ is a valid prefix of } \mathcal{G}_{\text{JSON}}(\sigma_t)\}$, yielding masked logits:
\begin{equation}
  \tilde{z}_s(v) =
  \begin{cases}
    z_s(v) & \text{if } v \in \mathcal{V}(w_{<s}), \\
    -\infty & \text{otherwise,}
  \end{cases}
  \qquad
  \pi_\theta(v \mid w_{<s}) = \frac{\exp(\tilde{z}_s(v))}{\sum_{v' \in \mathcal{V}(w_{<s})} \exp(\tilde{z}_s(v'))}.
  \label{eq:mcp:constrained_decoding}
\end{equation}
During post-training, tool-use competence is shaped in three stages \bookcite{reinforcement_learning_from_human_feedback}{PDF p.~175}:
\begin{enumerate}
  \item \textbf{SFT on Tool Trajectories:} Minimizes cross-entropy over synthetic/bootstrapped traces (e.g., Toolformer, ToolBench) to teach special-token delimiters and argument formatting.
  \item \textbf{Preference Optimization (DPO):} Optimizes pairwise preferences $(\tau^+ \succ \tau^-)$ to calibrate \emph{when} to invoke a tool versus answering directly from weights.
  \item \textbf{RL with Environment Feedback:} Maximizes expected terminal task reward $R(\tau)$ (e.g., constraint satisfaction, exact execution success) over multi-step tool trajectories:
  \begin{equation}
    \mathcal{J}(\theta) = \mathbb{E}_{\tau \sim \pi_\theta}\!\bigl[R(\tau)\bigr] - \beta \, D_{\mathrm{KL}}\!\bigl(\pi_\theta \,\|\, \pi_{\text{ref}}\bigr).
    \label{eq:mcp:rl_tool}
  \end{equation}
\end{enumerate}
\end{coremath}

\subsection{Hybrid Deterministic--Agentic DAG Reliability}

In workflow automation platforms such as n8n, workflows form directed acyclic graphs (DAGs) $G = (V, E)$ combining deterministic nodes $V_{\text{det}}$ (webhooks, HTTP requests, IF/Merge/Code transformations) with stochastic AI agent nodes $V_{\text{agent}}$ \bookcite{ultimate_n8n_starter_kit_2025}{PDF pp.~8--12, 26--31}. Because LLMs exhibit nonzero failure rates on exact counting, formatting, and constraint adherence \bookcite{ultimate_n8n_starter_kit_2025}{PDF p.~31}, wrapping an agent node with deterministic validation and retry/error branches directly governs end-to-end reliability.

\begin{coremath}[Reliability of Guarded Hybrid Workflow Nodes]
Suppose a workflow path traverses $L_{\text{det}}$ deterministic nodes (each with transient failure rate $\epsilon_{\text{det}} \ll 1$) and $L_{\text{agent}}$ stochastic AI agent nodes (each with single-shot schema/task failure probability $p_{\text{fail}}$). If each agent node is paired with a deterministic validator and up to $r$ bounded retries (or an error-branch fallback with success probability $q_{\text{fb}}$), the per-node failure probability drops from $p_{\text{fail}}$ to $p_{\text{fail}}^{r+1} (1 - q_{\text{fb}})$, and the end-to-end path success probability is:
\begin{equation}
  P_{\text{success}} = (1 - \epsilon_{\text{det}})^{L_{\text{det}}} \cdot \prod_{j=1}^{L_{\text{agent}}} \Bigl(1 - p_{\text{fail}, j}^{\,r_j + 1}\bigl(1 - q_{\text{fb}, j}\bigr)\Bigr).
  \label{eq:mcp:dag_reliability}
\end{equation}
\end{coremath}

% =============================================================================
\section{Algorithmic Implementation \& Comparative Synthesis}
\label{sec:mcp:algorithms}
% =============================================================================

\subsection{MCP JSON-RPC~2.0 Lifecycle, Transports \& Guarded Execution}

An MCP Host manages one or more embedded MCP Clients, each maintaining a 1:1 stateful or stateless session with an MCP Server over JSON-RPC~2.0 \bookcite{agent_tools_interoperability_mcp}{PDF pp.~21--24}. The protocol defines four message classes---Requests (expecting Results or Errors), Results, Errors, and one-way Notifications---carried over two standard transports:
\begin{itemize}
  \item \textbf{\texttt{stdio} (Local):} The Host spawns the MCP Server as a local subprocess and exchanges newline-delimited JSON-RPC messages over standard input/output, ideal for local filesystem or IDE integrations \bookcite{agent_tools_interoperability_mcp}{PDF p.~23}.
  \item \textbf{Streamable HTTP (Remote):} Uses HTTP \texttt{POST} as the command channel and optional Server-Sent Events (SSE) on \texttt{GET} as the announcement/streaming channel, supporting both stateful sessions and stateless serverless deployments \bookcite{agent_tools_interoperability_mcp}{PDF pp.~23--24}.
\end{itemize}

Algorithm~\ref{alg:mcp:guarded_loop} specifies the complete capability discovery and security-guarded execution loop, integrating manifest pinning, gateway allowlisting, tainted-input checks, and structured error feedback.

\begin{algorithm}[htbp]
\caption{MCP JSON-RPC Capability Discovery \& Guarded Tool Execution Loop}
\label{alg:mcp:guarded_loop}
\KwIn{User query $q$; MCP Server endpoints $\{\mathcal{S}_1, \dots, \mathcal{S}_m\}$; pinned manifest hashes $H_{\text{pin}}$; enterprise allowlist $\mathcal{A}_{\text{allow}}$; max steps $K_{\max}$}
\KwOut{Final grounded response $y_{\text{final}}$ or structured failure}
\ForEach{server $\mathcal{S}_j$}{
  Establish transport (\texttt{stdio} subprocess or \texttt{Streamable HTTP} session)\;
  Send JSON-RPC \texttt{initialize} request; negotiate protocol version \& capabilities\;
  Send \texttt{tools/list} request; receive manifests $\mathcal{M}_j = \{(\text{name}, \text{desc}, \text{inputSchema}, \text{outputSchema})\}$\;
  \If{$\mathrm{Hash}(\mathcal{M}_j) \neq H_{\text{pin}}[j]$ \textbf{or} unapproved tool in $\mathcal{M}_j \setminus \mathcal{A}_{\text{allow}}$}{
    Filter $\mathcal{M}_j \leftarrow \mathcal{M}_j \cap \mathcal{A}_{\text{allow}}$ via API/Agent Gateway and emit security alert\;
  }
  Sanitize tool descriptions against prompt-injection triggers before context injection\;
}
Initialize context $x_0 \leftarrow (\text{system\_prompt}, \{\mathcal{M}_j\}, q)$\;
\For{$k \leftarrow 1$ \KwTo $K_{\max}$}{
  Generate model output $(a_k \text{ or } y_{\text{final}}) \sim \pi_\theta(\cdot \mid x_{k-1})$ with schema-constrained decoding\;
  \If{$y_{\text{final}}$ emitted with no tool call}{
    \Return $y_{\text{final}}$\;
  }
  Validate $a_k = (\text{tool\_name}, \text{args})$ against $\text{inputSchema}$ and caller's scoped OAuth~2.0 audience token\;
  \If{tool has \texttt{destructiveHint=true} \textbf{and} tainted-data flow detected}{
    Require explicit user consent or abort with policy error\;
  }
  Dispatch JSON-RPC \texttt{tools/call} to target $\mathcal{S}_j$ (injecting credentials via side channel, never in LLM context)\;
  Receive JSON-RPC response $o_k$\;
  \eIf{$o_k.\text{isError} == \texttt{true}$}{
    Preserve actionable error string in $x_k$ so $\pi_\theta$ can self-correct arguments or fail over\;
  }{
    Validate $o_k$ against $\text{outputSchema}$; truncate/paginate if context budget exceeded; append to $x_k$\;
  }
}
\Return Timeout/fallback via deterministic workflow error branch\;
\end{algorithm}

\subsection{Tool-Design Best Practices \& Provider API Reconciliation}

Poorly documented tools degrade model selection accuracy and inflate token retries. To prevent these failures, \bookcite{agent_tools_interoperability_mcp}{PDF pp.~15--20} prescribes six core tool-design rules: (1) use specific, human-readable names (such as \texttt{create\_critical\_bug\_in\_jira} rather than \texttt{update\_jira}); (2) document every input and output parameter with types and defaults; (3) keep parameter lists short and granular; (4) describe \emph{what} action to achieve in system instructions rather than \emph{how} or duplicating tool docstrings; (5) publish task-oriented tools rather than raw 1:1 REST wrappers; and (6) return concise, actionable error strings (via \texttt{"isError": true}) that guide the LLM on recovery.

At the model-provider layer, \bookcite{reinforcement_learning_from_human_feedback}{PDF p.~174} notes that while the underlying JSON Schema semantics are shared, wire representations differ: OpenAI Chat Completions uses \texttt{tool\_calls} arrays (and the Responses API uses \texttt{function\_call} / \texttt{function\_call\_output} keyed by \texttt{call\_id}); Anthropic uses \texttt{input\_schema} with \texttt{tool\_use} / \texttt{tool\_result} content blocks; and Gemini exposes function-calling modes \texttt{AUTO}, \texttt{ANY}, \texttt{NONE}, and \texttt{VALIDATED}.

\subsection{A2A Interoperability, Agent Cards \& Enterprise Registries}

While MCP standardizes tool access, A2A (governed by the Linux Foundation) standardizes stateful, cross-framework agent collaboration \bookcite{prototype_production}{PDF pp.~27--35}. Every A2A-compliant agent publishes an \textbf{Agent Card}---a JSON manifest typically served at \texttt{/.well-known/agent-card.json}---declaring its \texttt{name}, \texttt{version}, \texttt{securitySchemes} (e.g., OAuth~2.0), \texttt{defaultInputModes}/\texttt{defaultOutputModes}, \texttt{skills}, and endpoint \texttt{url} \bookcite{prototype_production}{PDF pp.~28--29}. Frameworks such as Google's ADK expose existing agents via \texttt{to\_a2a(root\_agent, port=8001)} and consume remote peers via \texttt{RemoteA2aAgent} \bookcite{prototype_production}{PDF pp.~30--32}.

\begin{crossbook}[Cross-Book Synthesis: Protocol Roles, Registries \& Workflow Orchestration]
\label{tab:mcp:synthesis}
\vspace{0.3em}
\noindent\begin{tabularx}{\linewidth}{@{} >{\raggedright\arraybackslash}p{2.4cm} >{\raggedright\arraybackslash}p{2.5cm} >{\raggedright\arraybackslash}p{3.3cm} X @{}}
\toprule
\textbf{Layer} & \textbf{Protocol / Tool} & \textbf{Primary Control} & \textbf{Production Mandate \& Scale Pattern} \\
\midrule
Agent-to-Agent (Horizontal) &
  A2A (Agent Cards, Tasks) &
  Autonomous Peer Agents &
  Requires distributed trace IDs and persistent state management \bookcite{prototype_production}{PDF p.~32}. \\
\addlinespace
Agent-to-Tool (Vertical) &
  MCP (JSON-RPC~2.0) &
  Model (Tools) / App (Resources) / User (Prompts) &
  At $\sim$50 tools, static config suffices; at $\sim$5,000 tools, use a Tool Registry (Generalist, Specialist, or Dynamic subsets) \bookcite{prototype_production}{PDF p.~35}. \\
\addlinespace
Workflow Graph (Orchestration) &
  n8n DAG + AI Agent Nodes &
  Deterministic Triggers \& Error Branches &
  Use deterministic nodes for APIs/transforms/validation and Error Trigger workflows for graceful recovery \bookcite{ultimate_n8n_starter_kit_2025}{PDF pp.~22--25, 31}. \\
\bottomrule
\end{tabularx}
\end{crossbook}

% =============================================================================
\section{Visual Diagnostics \& Geometric Structure}
\label{sec:mcp:visuals}
% =============================================================================

Figure~\ref{fig:mcp:host_client_server} illustrates the internal MCP Host--Client--Server session topology and JSON-RPC~2.0 \texttt{tools/call} payload structure, while Figure~\ref{fig:mcp:transports} contrasts local \texttt{stdio} subprocesses against remote \texttt{HTTP+SSE} and \texttt{Streamable HTTP} transports.

\monographfig[0.80\textwidth]{agent_tools_interoperability_mcp_p22_fig1.png}{MCP Host, Client, and Server architecture within an AI Agent application, illustrating 1:1 Client--Server sessions and the JSON-RPC~2.0 \texttt{tools/call} invocation format \bookcite{agent_tools_interoperability_mcp}{Fig.~2, PDF p.~22}.}{fig:mcp:host_client_server}

\monographfig[0.82\textwidth]{agent_tools_interoperability_mcp_p24_fig1.png}{MCP Transport Protocols: local \texttt{stdio} subprocess communication versus remote \texttt{HTTP + SSE} and \texttt{Streamable HTTP} (separating the \texttt{POST} command channel from the \texttt{GET} announcement channel for stateful and stateless servers) \bookcite{agent_tools_interoperability_mcp}{Fig.~3, PDF p.~24}.}{fig:mcp:transports}

When scaling beyond a single agent, A2A and MCP form a two-tier interoperability stack. Figure~\ref{fig:mcp:a2a_mcp_collab} shows the enterprise reference pattern from \bookcite{prototype_production}{PDF p.~33}, and Figure~\ref{fig:mcp:mcp_a2a_topology} formalizes the two-tier protocol topology separating horizontal A2A delegation from vertical MCP tool/resource plumbing.

\monographfig[0.82\textwidth]{prototype_production_p33_fig1.png}{Two-tier A2A and MCP collaboration: a Client/Router Agent delegates tasks horizontally via A2A across heterogeneous frameworks (CrewAI, LangGraph, ADK), while each specialized agent accesses tools vertically via MCP Servers and API Hubs \bookcite{prototype_production}{Fig.~4, PDF p.~33}.}{fig:mcp:a2a_mcp_collab}

\begin{figure}[htbp]
  \centering
  \resizebox{\linewidth}{!}{%
  \begin{tikzpicture}[
    node distance=0.85cm and 1.1cm,
    tier/.style={draw=monographRule, dashed, rounded corners=4pt, inner sep=8pt},
  ]
    % Top tier: User / Router Agent
    \node[monoBox, minimum width=3.4cm] (router) {Orchestrator / Router Agent\\[-2pt]{\scriptsize\color{monographSlate}(A2A Client + Trace ID)}};

    % Middle tier: Specialized A2A Agents
    \node[monoIndigoBox, below left=1.15cm and 0.6cm of router, minimum width=3.1cm] (agentA) {Specialized Agent A\\[-2pt]{\scriptsize\color{monographSlate}Agent Card + OAuth~2.0}};
    \node[monoIndigoBox, below right=1.15cm and 0.6cm of router, minimum width=3.1cm] (agentB) {Specialized Agent B\\[-2pt]{\scriptsize\color{monographSlate}Agent Card + State Store}};

    % Bottom tier: MCP Servers
    \node[monoTealBox, below=1.25cm of agentA, xshift=-1.3cm, minimum width=2.5cm] (mcp1) {MCP Server 1\\[-2pt]{\scriptsize\color{monographSlate}\texttt{stdio} / Local FS}};
    \node[monoTealBox, below=1.25cm of agentA, xshift=1.5cm, minimum width=2.5cm] (mcp2) {MCP Server 2\\[-2pt]{\scriptsize\color{monographSlate}Streamable HTTP}};
    \node[monoTealBox, below=1.25cm of agentB, xshift=1.3cm, minimum width=2.5cm] (mcp3) {MCP Server 3\\[-2pt]{\scriptsize\color{monographSlate}Resources \& Prompts}};

    % A2A horizontal/delegation arrows
    \draw[monoArrow] (router.south west) -- node[left, font=\scriptsize, xshift=-2pt] {A2A Task JSON-RPC} (agentA.north);
    \draw[monoArrow] (router.south east) -- node[right, font=\scriptsize, xshift=2pt] {A2A Task JSON-RPC} (agentB.north);
    \draw[<->, >=stealth, thick, monographIndigo, dashed] (agentA.east) -- node[above, font=\scriptsize] {A2A Peer Delegation} (agentB.west);

    % MCP vertical arrows
    \draw[monoArrow, monographTeal] (agentA.south) -- (mcp1.north);
    \draw[monoArrow, monographTeal] (agentA.south) -- node[right, font=\scriptsize, pos=0.55] {MCP \texttt{tools/call}} (mcp2.north);
    \draw[monoArrow, monographTeal] (agentB.south) -- (mcp2.north);
    \draw[monoArrow, monographTeal] (agentB.south) -- node[right, font=\scriptsize, pos=0.55] {MCP \texttt{resources/get}} (mcp3.north);

    % Tier annotations on right margin
    \node[font=\scriptsize\bfseries\color{monographIndigo}, right=0.5cm of agentB, align=left] {Horizontal Tier:\\Stateful Goal Delegation\\(A2A)};
    \node[font=\scriptsize\bfseries\color{monographTeal}, right=0.5cm of mcp3, align=left] {Vertical Tier:\\$\mathcal{O}(M+N)$ Tool Plumbing\\(MCP)};
  \end{tikzpicture}%
  }
  \caption{Unified two-tier interoperability topology: A2A governs horizontal, stateful goal delegation between heterogeneous agents via Agent Cards, while MCP governs vertical, schema-constrained access to Tools, Resources, and Prompts.}
  \label{fig:mcp:mcp_a2a_topology}
\end{figure}

Figure~\ref{fig:mcp:integration_scaling} quantifies the connector scaling advantage of standardized MCP/A2A hubs (Eq.~\ref{eq:mcp:m_plus_n_cost}) over bespoke $M \times N$ glue code (Eq.~\ref{eq:mcp:mn_cost}) as the ecosystem grows.

\begin{figure}[htbp]
  \centering
  \begin{tikzpicture}
    \begin{axis}[
      width=0.84\textwidth,
      height=6.0cm,
      xlabel={Ecosystem Scale $N$ (Tools / Services)},
      ylabel={Required Integrations},
      xmin=2, xmax=30,
      ymin=0, ymax=920,
      grid=major,
      grid style={dashed, gray!30},
      legend pos=north west,
      legend style={font=\scriptsize, draw=monographRule, fill=white, fill opacity=0.9},
      tick label style={font=\scriptsize},
      label style={font=\small\bfseries},
    ]
      % M = N bespoke: N^2
      \addplot[color=monographCrimson, very thick, domain=2:30, samples=30] {x^2};
      \addlegendentry{Bespoke Glue ($M = N$): $M \cdot N = N^2$}

      % M = 10 fixed, varying N bespoke: 10*N
      \addplot[color=monographIndigo, thick, dashed, domain=2:30, samples=30] {10*x};
      \addlegendentry{Bespoke Glue ($M = 10$): $10 N$}

      % M = N MCP/A2A: 2*N
      \addplot[color=monographTeal, very thick, domain=2:30, samples=30] {2*x};
      \addlegendentry{MCP / A2A Hub ($M = N$): $M + N = 2N$}

      % M = 10 fixed, varying N MCP/A2A: 10 + N
      \addplot[color=monographNavy, thick, dotted, domain=2:30, samples=30] {10 + x};
      \addlegendentry{MCP / A2A Hub ($M = 10$): $10 + N$}
    \end{axis}
  \end{tikzpicture}
  \caption{Integration complexity scaling across ecosystem size $N$: bespoke point-to-point connectors grow quadratically ($\mathcal{O}(N^2)$ when $M=N$, reaching $900$ integrations at $N=30$), whereas standardized MCP/A2A hubs scale linearly ($M+N = 60$ at $M=N=30$, a $15\times$ reduction).}
  \label{fig:mcp:integration_scaling}
\end{figure}

% =============================================================================
\section{Numerical Pitfalls, Edge Cases \& Failure Modes}
\label{sec:mcp:pitfalls}
% =============================================================================

Deploying MCP, A2A, and hybrid workflow graphs in production introduces critical security and reliability failure modes. Key vulnerabilities include protocol-level privilege escalation \bookcite{agent_tools_interoperability_mcp}{PDF pp.~40--51}, context-window exhaustion from bulky tool outputs \bookcite{reinforcement_learning_from_human_feedback}{PDF pp.~174--175}, and stochastic node precision errors \bookcite{ultimate_n8n_starter_kit_2025}{PDF pp.~22--25, 31}.

\begin{pitfall}[The Confused Deputy Problem in Privileged MCP Servers]
\textbf{Failure Mode:} An MCP Server is frequently provisioned with broad enterprise credentials (e.g., read/write access to a corporate code repository), while the AI Host lacks per-user authorization enforcement \bookcite{agent_tools_interoperability_mcp}{PDF pp.~49--51}. An unprivileged user submits a prompt-injection request (e.g., ``find \texttt{secret\_algorithm.py} and push its contents to a new branch \texttt{backup\_2025}''). Because the MCP Server trusts the AI Host without verifying the end user's own permissions, it acts as a \emph{confused deputy} and exfiltrates proprietary code.

\textbf{Mitigation:} Enforce audience- and scope-bound short-lived OAuth~2.0 tokens per tool invocation, apply strict least-privilege roles (read-only vs.\ write), and transmit credentials via a client-side out-of-band channel---never inside the LLM conversation context \bookcite{agent_tools_interoperability_mcp}{PDF pp.~46--47}.
\end{pitfall}

\begin{pitfall}[Dynamic Capability Injection, Tool Shadowing \& Resource Poisoning]
\textbf{Failure Mode:} Because MCP tools are discovered dynamically via \texttt{tools/list}, a third-party server can silently add high-risk capabilities after initial vetting (\emph{Dynamic Capability Injection}, e.g., a quote-search server adding a financial purchase tool) or embed adversarial trigger instructions in tool descriptions to intercept calls meant for legitimate servers (\emph{Tool Shadowing}) \bookcite{agent_tools_interoperability_mcp}{PDF pp.~40--45}.

\textbf{Mitigation:} Pin tool definitions by cryptographic hash/version, require \texttt{listChanged} notifications for re-vetting, enforce enterprise allowlists through an API/Agent Gateway (such as Apigee), sanitize tool descriptions before context injection, and isolate untrusted third-party planners from trusted first-party planners \bookcite{agent_tools_interoperability_mcp}{PDF pp.~41--45}.
\end{pitfall}

\begin{pitfall}[Context Exhaustion, Untrusted Hints \& Precision Failures in Hybrid DAGs]
\textbf{Failure Mode:} (1) Search or retrieval tools can flood the context window, degrading reasoning accuracy \bookcite{reinforcement_learning_from_human_feedback}{PDF p.~175}. (2) MCP tool \texttt{annotations} (\texttt{readOnlyHint}, \texttt{destructiveHint}) are unverified server hints and must never be trusted as security guarantees \bookcite{agent_tools_interoperability_mcp}{PDF pp.~26--27}. (3) LLM nodes inside n8n workflows fail on exact counting or deterministic formatting \bookcite{ultimate_n8n_starter_kit_2025}{PDF p.~31}.

\textbf{Mitigation:} Paginate or summarize bulky tool outputs, treat all external inputs/outputs as tainted until schema-validated, and wrap stochastic AI Agent nodes with deterministic n8n validation nodes, error branches, and dedicated Error Trigger workflows \bookcite{ultimate_n8n_starter_kit_2025}{PDF pp.~22--25}.
\end{pitfall}

\end{document}
